%!TEX root = ../../../main.tex
\chapter{Gefährdung von Lebensdauer und Funktionsfähigkeit}\label{chap:cloud-folgen}
% Kennzeichnung der Autorenschaft (Vorgabe der Arbeit) -- Form ggf. in der Gruppe vereinheitlichen
\noindent\textit{Verfasser: [Name]}

TEST

Dieses Kapitel untersucht die Folgen der in Kapitel~\ref{chap:cloud-ursachen}
beschriebenen Abhängigkeit. Auf eine begriffliche Einordnung folgen die
Analyse des Extremfalls Gigaset und ein Vergleich mit weiteren Abschaltungen.

\section{Softwarebedingte Obsoleszenz und prekäres Eigentum}\label{sec:folgen-rahmen}

Die Studie für das \acrshort{UBA} versteht unter Softwareobsoleszenz die
Verkürzung der Nutzungsdauer eines technisch funktionsfähigen Geräts aufgrund
von Software \autocite[vgl.][S.~5]{jaegererben:2023}. Neben Ursachen in der
Software selbst nennt sie Veränderungen ihres Umfelds, etwa die Einstellung
des Supports \autocite[vgl.][S.~32\,f.]{jaegererben:2023}, und behandelt
externe Abhängigkeiten wie Cloud-Server als eigene Kategorie
\autocite[vgl.][S.~151\,ff.]{jaegererben:2023}. Die Ursachen seien dabei
weniger technischer als wirtschaftlicher und organisatorischer Natur
\autocite[vgl.][S.~5]{jaegererben:2023}. Hinzu kommen einseitige Änderungen
über Updates und Nutzungsbedingungen, wie die in Kapitel~\ref{chap:cloud-ursachen}
beschriebenen Schnittstellensperren und nachträglichen Gebühren.

Für Käufer stellt sich damit die Frage, was sie eigentlich erwerben.
Perzanowski und Schultz zeigen, dass Lizenzbedingungen und digitale
Kontrollmechanismen klassische Eigentumsrechte aushöhlen
\autocite[vgl.][]{perzanowski:2016}. Tusikov spricht von einem hybriden und
prekären Eigentum: Die Hardware gehört dem Käufer, die Software darf er nur
im Rahmen einer Lizenz nutzen \autocite[vgl.][S.~11]{tusikov:2019}. Dass
Hersteller diese Verfügungsgewalt ausüben, zeigte bereits 2016 der Fall
Revolv: Das zum Alphabet-Konzern gehörende Unternehmen Nest hatte den
Hersteller dieser Smart-Home-Zentrale übernommen und schaltete die Geräte zum
15.~Mai~2016 vollständig ab, um sich auf die eigene Plattform zu
konzentrieren \autocites[vgl.][]{protalinski:2016}{cullen:2016}. Geräte
werden also nicht nur durch Insolvenzen unbrauchbar, sondern auch durch
strategische Entscheidungen wirtschaftlich starker Unternehmen.

\section{Fallanalyse: Gigaset Smart Home}\label{sec:folgen-gigaset}

Am deutlichsten zeigen sich die Folgen bei Gigaset. Über das Vermögen der
Gigaset Communications GmbH wurde am 29.~Januar~2024 das Insolvenzverfahren
eröffnet. Für das Telefoniegeschäft fand sich ein Käufer, für den Bereich
Smart Home/Care nicht. Das Unternehmen stellte die Cloud-Dienste deshalb zum
29.~März~2024 ein, da ein Weiterbetrieb insolvenzrechtlich nicht möglich sei
\autocite[vgl.][]{gigaset:2024b}. Zwischen Verfahrenseröffnung und
Abschaltung lagen nur zwei Monate.

Nach Angaben des Herstellers waren danach weder die Apps noch die vernetzten
Sensoren und Geräte nutzbar. Gebuchte Kamera-Servicepakete endeten, eine
Rückgabe der Geräte war ausgeschlossen, und nur der lokale Alarm des
Rauchmelders funktionierte weiter \autocite[vgl.][]{gigaset:2024a}. Betroffen
waren vor allem Tür-, Fenster- und Bewegungssensoren, Alarmsysteme und
Kameras \autocite[vgl.][]{wittkowski:2024}. Der Ausfall traf damit den Schutz
der Wohnung, für den die Geräte angeschafft worden waren. Ansprüche konnten
die Kunden nur als Insolvenzforderungen anmelden, deren Erfüllung von der
verbleibenden Masse abhängt \autocite[vgl.][]{teltarif:2024}.

Gemessen an Kapitel~2 zeigt der Fall einen vollständigen Verlust der
Anwendungssouveränität, denn die Nutzbarkeit hing allein am Fortbestand eines
Unternehmens. Die Entscheidungssouveränität war schon beim Kauf
eingeschränkt, weil diese Abhängigkeit kaum erkennbar war
(vgl. Abschnitt~\ref{sec:cloud-haas}). Ökologisch steht der Fall für das
zusätzliche Elektronikschrottaufkommen, das die \acrshort{UBA}-Studie als
Folge softwarebedingter Obsoleszenz beschreibt
\autocite[vgl.][S.~18]{jaegererben:2023}.

\section{Fallvergleich: Schutzfaktoren gegen den Funktionsverlust}\label{sec:folgen-vergleich}

Nicht jede Abschaltung verläuft so drastisch.
Tabelle~\ref{tab:folgen-abschaltungen} vergleicht weitere Fälle aus den
Jahren 2022 bis 2026, geordnet vom schwersten zum mildesten Ausgang.

\begin{table}[htbp]
  \centering
  \caption{Vergleich von Abschaltungen im Smart Home}
  \label{tab:folgen-abschaltungen}
  \footnotesize
  \setlength{\tabcolsep}{4pt}
  \renewcommand{\arraystretch}{1.2}
  \begin{tabularx}{\textwidth}{@{}>{\raggedright\arraybackslash}p{2.45cm}
      >{\raggedright\arraybackslash}p{1.85cm}
      >{\raggedright\arraybackslash}p{2.6cm}
      >{\raggedright\arraybackslash}X
      >{\raggedright\arraybackslash}p{2.35cm}@{}}
    \hline
    \textbf{Fall} & \textbf{Abschal\-tung} & \textbf{Anlass}
      & \textbf{Verbleibende Funktion} & \textbf{Schutzfaktor} \\
    \hline
    Gigaset Smart Home & 29.03.2024 & Insolvenz
      & Lediglich lokaler Rauchmelderalarm & Keiner \\
    Belkin Wemo & 31.01.2026 & Umschichtung von Ressourcen
      & Nur Geräte mit Apple HomeKit oder Thread & Offener Standard \\
    Insteon & April 2022 & Geschäfts\-aufgabe ohne Vorwarnung
      & Lokale Steuerung über Home Assistant; Dienst später von Nutzern
        übernommen & Lokale Grundfunktion \\
    devolo Home Control & 31.12.2025 & Einstellung des Systems
      & Eingeschränkter lokaler Betrieb; Z-Wave-Geräte an fremder Zentrale
      & Offenes Funkprotokoll \\
    Google Nest (1./2.~Gen.) & 25.10.2025 & Veraltete Hardware
      & Bedienung am Gerät, gespeicherte Zeitpläne & Lokale Grundfunktion \\
    Logitech Squeezebox & 2024 & Einstellung des Onlinedienstes
      & Betrieb mit lokalem Open-Source-Server & Offene Software \\
    Bose SoundTouch & 06.05.2026 & Supportende
      & Lokale App, offengelegte \acrshort{API}, Community-Lösungen
      & Offene Schnittstelle \\
    \hline
  \end{tabularx}
  \par\smallskip
  \raggedright\scriptsize
  Quelle: Eigene Darstellung auf Basis von
  \cite{gigaset:2024a,belkin:2026,rogerson:2022a,schoutsen:2022,giga:2025,google:2025,golem:2024,rudra:2026,stereoguide:2026}.
\end{table}

Bei Belkin funktionieren nur Wemo-Geräte weiter, die zuvor in Apple HomeKit
eingebunden waren oder auf dem Standard Thread basieren
\autocite[vgl.][]{belkin:2026}. Aufschlussreich ist Insteon: Das Unternehmen
schaltete seine Server im April~2022 ohne Vorwarnung ab
\autocite[vgl.][]{rogerson:2022a}. Weil die Geräte lokal kommunizierten,
ließen sie sich jedoch über Home Assistant ohne Cloud weiter steuern
\autocite[vgl.][]{schoutsen:2022}, bis eine Gruppe von Nutzern das
Unternehmen im Juni~2022 übernahm \autocite[vgl.][]{rogerson:2022b}. Die
ältesten Nest-Thermostate behalten nach dem Supportende ihre Bedienung am
Gerät und gespeicherte Zeitpläne \autocites[vgl.][]{google:2025}{heise:2025a}.
Logitech und Bose schließlich zeigen den Wert offener Software: Die
Squeezebox-Player laufen mit einem von der Community gepflegten
Open-Source-Server weiter \autocites[vgl.][]{golem:2024}{lyrion:oj}, und auf
Basis der von Bose veröffentlichten Schnittstellendokumentation bilden
Community-Projekte die SoundTouch-Cloud lokal nach
\autocites[vgl.][]{rudra:2026}{gesellix:2026}.

Aus dem Vergleich lassen sich drei Schutzfaktoren ableiten:
\begin{enumerate}
  \item eine \textbf{lokale Grundfunktion}, die ohne Server erhalten bleibt,
  \item ein \textbf{offener Standard oder ein offenes Funkprotokoll}, über
        das sich das Gerät in andere Systeme einbinden lässt,
  \item eine \textbf{offengelegte Schnittstelle oder offene Software}, die
        Dritten die Weiterentwicklung erlaubt.
\end{enumerate}
Fehlen alle drei wie bei Gigaset, wird das Gerät mit der Abschaltung wertlos.
Die Faktoren entsprechen den Stoßrichtungen der \acrshort{UBA}-Studie, die
langfristige Interoperabilität und die Gewährleistung definierter
Kernfunktionen als zentrale Ansatzpunkte nennt
\autocite[vgl.][S.~14]{jaegererben:2023}.

Weiterbetrieb bedeutet allerdings nicht Sicherheit. Mit dem Supportende
entfallen meist auch Sicherheitsupdates; für SoundTouch schließt Bose weitere
Patches ausdrücklich aus \autocite[vgl.][]{stereoguide:2026}. Eine
Langzeituntersuchung zeigt, dass Millionen vernetzter Geräte nach ihrem
Supportende in Betrieb bleiben und mehr als die Hälfte der gefundenen
Schwachstellen erst danach entdeckt wurde \autocite[vgl.][]{wang:2023}. Die
Schutzfaktoren sichern also die Funktion, nicht automatisch die Sicherheit.
Wie sich solche Geräte im Heimnetz absichern lassen, behandelt Kapitel~8.

\section{Zwischenfazit}\label{sec:folgen-zwischenfazit}

Die Cloud-Abhängigkeit vernetzter Haushaltsgeräte berührt alle drei
Souveränitätsdimensionen: die \textbf{Datensouveränität}, weil Wohnungskarten
bei einem Eigentümerwechsel ohne Mitsprache der Nutzer in neue Hände
gelangen; die \textbf{Entscheidungssouveränität}, weil Hersteller Funktionen
nach dem Kauf einseitig sperren oder bepreisen; und die
\textbf{Anwendungssouveränität}, weil die Nutzbarkeit im Extremfall mit dem
Hersteller endet.

Das geltende Recht schließt diese Lücke nur teilweise. Die
Aktualisierungspflicht für Waren mit digitalen Elementen trifft den
Verkäufer und gilt für den Zeitraum, den Verbraucher nach Art und Zweck der
Ware erwarten können \autocite[vgl.][§~475b Abs.~4 Nr.~2]{bgb:2026}; wie lang
dieser ist, lässt das Gesetz offen. Der Data Act gewährt Zugang zu den Daten
vernetzter Produkte \autocite[vgl.][Art.~4]{dataact:2023}, verpflichtet aber
nicht zum Weiterbetrieb von Diensten, und in der Insolvenz bleiben ohnehin
nur Insolvenzforderungen.

Daraus folgen vier Anforderungen: eine lokale Rückfallebene für
Kernfunktionen, offene Standards und Funkprotokolle, offengelegte
Schnittstellen sowie Transparenz über die Dauer der Unterstützung
einschließlich der Sicherheitsupdates. In diesem Sinne empfiehlt die
\acrshort{UBA}-Studie, den Betrieb ohne externe Abhängigkeiten zu ermöglichen,
externe Dienste mindestens zehn Jahre aufrechtzuerhalten und den
Supportzeitraum verpflichtend anzugeben
\autocite[vgl.][S.~18\,f.]{jaegererben:2023}. Die Kapitel~6 und~7 prüfen,
inwieweit lokale Zentralen, offene Standards und alternative Firmware diese
Anforderungen erfüllen.
